%-----------------------------------------------------------------------------------------------------------------------------------------------%
%	The MIT License (MIT)
%
%	Copyright (c) 2021 Jitin Nair
%
%	Permission is hereby granted, free of charge, to any person obtaining a copy
%	of this software and associated documentation files (the "Software"), to deal
%	in the Software without restriction, including without limitation the rights
%	to use, copy, modify, merge, publish, distribute, sublicense, and/or sell
%	copies of the Software, and to permit persons to whom the Software is
%	furnished to do so, subject to the following conditions:
%	
%	THE SOFTWARE IS PROVIDED "AS IS", WITHOUT WARRANTY OF ANY KIND, EXPRESS OR
%	IMPLIED, INCLUDING BUT NOT LIMITED TO THE WARRANTIES OF MERCHANTABILITY,
%	FITNESS FOR A PARTICULAR PURPOSE AND NONINFRINGEMENT. IN NO EVENT SHALL THE
%	AUTHORS OR COPYRIGHT HOLDERS BE LIABLE FOR ANY CLAIM, DAMAGES OR OTHER
%	LIABILITY, WHETHER IN AN ACTION OF CONTRACT, TORT OR OTHERWISE, ARISING FROM,
%	OUT OF OR IN CONNECTION WITH THE SOFTWARE OR THE USE OR OTHER DEALINGS IN
%	THE SOFTWARE.
%	
%
%-----------------------------------------------------------------------------------------------------------------------------------------------%

%----------------------------------------------------------------------------------------
%	DOCUMENT DEFINITION
%----------------------------------------------------------------------------------------

% article class because we want to fully customize the page and not use a cv template
\documentclass[a4paper,12pt]{article}

%----------------------------------------------------------------------------------------
%	FONT
%----------------------------------------------------------------------------------------

% % fontspec allows you to use TTF/OTF fonts directly
% \usepackage{fontspec}
% \defaultfontfeatures{Ligatures=TeX}

% % modified for ShareLaTeX use
% \setmainfont[
% SmallCapsFont = Fontin-SmallCaps.otf,
% BoldFont = Fontin-Bold.otf,
% ItalicFont = Fontin-Italic.otf
% ]
% {Fontin.otf}

%----------------------------------------------------------------------------------------
%	PACKAGES
%----------------------------------------------------------------------------------------
\usepackage{url}
\usepackage{parskip} 	

%other packages for formatting
\RequirePackage{color}
\RequirePackage{graphicx}
\usepackage[usenames,dvipsnames]{xcolor}
\usepackage[scale=0.9]{geometry}

%tabularx environment
\usepackage{tabularx}

%for lists within experience section
\usepackage{enumitem}

% centered version of 'X' col. type
\newcolumntype{C}{>{\centering\arraybackslash}X} 

%to prevent spillover of tabular into next pages
\usepackage{supertabular}
\usepackage{tabularx}
\newlength{\fullcollw}
\setlength{\fullcollw}{0.47\textwidth}

%custom \section
\usepackage{titlesec}				
\usepackage{multicol}
\usepackage{multirow}

%cyrillic fonts

\usepackage[utf8]{inputenc}
\usepackage[russian]{babel}

%CV Sections inspired by: 
%http://stefano.italians.nl/archives/26
\titleformat{\section}{\Large\scshape\raggedright}{}{0em}{}[\titlerule]
\titlespacing{\section}{0pt}{10pt}{6pt}

%for publications
\usepackage[style=authoryear,sorting=ynt, maxbibnames=2]{biblatex}

%Setup hyperref package, and colours for links
\usepackage[unicode, draft=false]{hyperref}
\definecolor{linkcolour}{rgb}{0,0.2,0.6}
\hypersetup{colorlinks,breaklinks,urlcolor=linkcolour,linkcolor=linkcolour}
\addbibresource{citations.bib}
\setlength\bibitemsep{1em}

%for social icons
\usepackage{fontawesome5}

%debug page outer frames
%\usepackage{showframe}

%----------------------------------------------------------------------------------------
%	BEGIN DOCUMENT
%----------------------------------------------------------------------------------------
\begin{document}

% non-numbered pages
\pagestyle{empty} 

%----------------------------------------------------------------------------------------
%	TITLE
%----------------------------------------------------------------------------------------

% \begin{tabularx}{\linewidth}{ @{}X X@{} }
% \huge{{Хабаров Роман}}\vspace{2pt} & \hfill \emoji{incoming-envelope} xabarov.roman@gmail.com \\
% \raisebox{-0.05\height}\faGithub\ username \ | \
% \raisebox{-0.00\height}\faLinkedin\ username \ | \ \raisebox{-0.05\height}\faGlobe \ mysite.com  & \hfill \emoji{calling} number
% \end{tabularx}

\begin{tabularx}{\linewidth}{@{} C @{}}
\Huge{Хабаров Роман} \\[7.5pt]
\Large{Data Scientist} \\[2.5pt]
\href{https://github.com/xabarov}{\raisebox{-0.05\height}\faGithub\ xabarov} \ $|$ \ 
% \href{https://linkedin.com/in/username}{\raisebox{-0.05\height}\faLinkedin\ username} \ $|$ \ 
% \href{https://mysite.com}{\raisebox{-0.05\height}\faGlobe \ mysite.com} \ $|$ \ 
\href{mailto:xabarov1985@gmail.com}{\raisebox{-0.05\height}\faEnvelope \ xabarov1985@gmail.com} \ $|$ \ 
\href{https://elibrary.ru/author_items.asp?authorid=1045982}{\raisebox{-0.05\height}\faFile\ elibrary.ru} \ $|$ \ 
\href{tel:+79118345456}{\raisebox{-0.05\height}\faMobile \ +7 911 834 54 56} \\
\end{tabularx}

%----------------------------------------------------------------------------------------
% EXPERIENCE SECTIONS
%----------------------------------------------------------------------------------------

%Interests/ Keywords/ Summary
\section{Краткое резюме}
Кандидат технических наук. 7 лет в науке (научный сотрудник, начальник лаборатории в ИТМО, МИРЭА, ВКА им. А. Ф. Можайского). После защиты диссертации (2020) занялся ML; стаж в ML — 5 лет. Опыт в CV (детекция объектов, семантическая сегментация на космоснимках) и NLP (классификация, NER, RAG, LLM-агенты). 

%Experience
\section{Опыт работы}

\begin{tabularx}{\linewidth}{ @{}l r@{} }
\textbf{ВКА им. А. Ф. Можайского} & \hfill Сентябрь 2017 - Февраль 2024 \\[3.75pt]
\multicolumn{2}{@{}X@{}}{Работал на должностях научного сотрудника, начальника лаборатории над проектами в областях CV и NLP.
 
\begin{description}
\item[$\bullet$] \textbf{CV (космическая съёмка):} разработал Python-модуль для детекции и семантической сегментации промышленных объектов на спутниковых снимках; настроил pipeline сбора данных, разметки и обучения. Модели: YOLO, Cascade R-CNN, DeepLabv3, PSANet, PSPNet, UNet, SegFormer.
\item[$\bullet$] \textbf{GraphRAG:} реализовал модуль построения семантического графа знаний и связывания с онтологией предметной области для улучшения качества ответов RAG-систем.
\item[$\bullet$] \textbf{Автоматизация пентеста:} разработал консольное приложение на основе LLM-агентов для автоматизации сценариев тестирования на проникновение.
\item[$\bullet$] \textbf{Система сбора и анализа документов:} построил pipeline сбора, индексации и семантического поиска по корпусу документов.
\item[$\bullet$] \textbf{Доверенный ИИ:} участвовал в исследовании уязвимостей ML-моделей — проводил тестирование атак на классификаторы текста и модели компьютерного зрения. 

\end{description}
}
\end{tabularx}

\begin{tabularx}{\linewidth}{ @{}l r@{} }
\textbf{ИТМО, МИРЭА - Российский технический университет} & \hfill Февраль 2021 - Февраль 2024 \\[3.75pt]
\multicolumn{2}{@{}X@{}}{
Работал на должностях младшего и старшего научного сотрудника. Написал 6 модулей для специализированных методик расчёта на C++ (Qt), в том числе с параллезацией на GPU (OpenCL); опыт математического моделирования, разработки под заказ. Участвовал в проектах:
\begin{minipage}[t]{\linewidth}
\begin{itemize}[nosep,after=\strut, leftmargin=1em, itemsep=3pt]
\item[--] Создание программно-аппаратного комплекса проведения математического моделирования взаимодействия бортовой радиоэлектронной аппаратуры составных частей КРК «Ангара».
\item[--] СЧ ОКР «Разработка комплексной системы расчетно-экспериментальной оценки стойкости РН «Ангара-А5М» к попаданию молниевого разряда».
\end{itemize}
\end{minipage}
}
\end{tabularx}

%Projects
\section{Проекты}

\begin{tabularx}{\linewidth}{ @{}l r@{} }
\textbf{Most Queue} & \hfill \href{https://github.com/xabarov/most-queue}{GitHub} \\[3.75pt]
\multicolumn{2}{@{}X@{}}{Пакет для имитационного моделирования и расчёта систем массового обслуживания (Python).}  \\

\textbf{OSINT GraphRAG} & \hfill \href{https://github.com/xabarov/osint-graphrag}{GitHub} \\[3.75pt]
\multicolumn{2}{@{}X@{}}{Система извлечения графов знаний из OSINT-текстов: сущности (персоны, организации), связи, дедупликация; семантический поиск (Milvus), чат-интерфейс для запросов на естественном языке (FastAPI, Neo4j, React).}  \\

\textbf{Deep Research Agent} & \hfill \href{https://github.com/xabarov/deep_research}{GitHub} \\[3.75pt]
\multicolumn{2}{@{}X@{}}{Агент глубокого исследования на smolagents: веб-поиск, научные базы (arXiv, PubMed), соцсети, Wikipedia; многоуровневая суммаризация контекста, факт-ориентированная память и векторный поиск (Qdrant); режимы исследования и OSINT-досье (FastAPI, React).}  \\
\end{tabularx}

%----------------------------------------------------------------------------------------
%	EDUCATION
%----------------------------------------------------------------------------------------
\section{Образование}
\begin{tabularx}{\linewidth}{@{}l X@{}}	
2017 - 2020 & PhD (Computer Science, системный анализ, моделирование, компьютерные технологии) \textbf{ВКА им. А. Ф. Можайского} \\

2002 - 2007 & Инженер РЭО \textbf{ИВВАИУ} 
\end{tabularx}

%----------------------------------------------------------------------------------------
%	PUBLICATIONS
%----------------------------------------------------------------------------------------
\section{Публикации}
\begin{tabularx}{\linewidth}{@{}l X@{}}	
\multicolumn{2}{@{}X@{}}{
Большинство публикаций посвящено тематике диссертации (теория очередей). Ведётся работа над публикациями в области ML.
\begin{minipage}[t]{\linewidth}
\begin{itemize}[nosep,after=\strut, leftmargin=1em, itemsep=3pt]
\item[1.] Лохвицкий, В. А. Численный расчет многоканальной системы массового обслуживания с разогревом, охлаждением и задержкой начала охлаждения / В. А. Лохвицкий, Р. С. Хабаров, Е. Л. Яковлев // Авиакосмическое приборостроение. – 2025. – № 1. – С. 44-57. – DOI 10.25791/aviakosmos.1.2025.1456. – EDN OVJXKE.
\item[2.] Хабаров, Р. С. Аппроксимация вероятностно-временных характеристик сложных систем массового обслуживания на основе регрессионных моделей машинного обучения / Р. С. Хабаров, Э. С. Левчик, В. А. Лохвицкий // Имитационное моделирование. Теория и практика (ИММОД-2023) : Сборник трудов одиннадцатой всероссийской научно-практической конференции по имитационному моделированию и его применению в науке и промышленности, Казань, 18–20 октября 2023 года. – Казань: Издательство АН РТ, 2023. – С. 735-742. – EDN TMPERE  
\item[3.] Mochalov, V. F. Processing of Multispectral Survey Materials based on Machine Learning Methods for Forest Management / V. F. Mochalov, R. S. Khabarov // Proceedings of 2021 IV International Conference on Control in Technical Systems (CTS), Saint Petersburg Electrotechnical University “LETI”, 21–23 сентября 2021 года. – IEEE, 2021. – P. 217-219. – DOI 10.1109/CTS53513.2021.9562930. – EDN WWLBGW.  
\end{itemize}
\end{minipage}
}  
\end{tabularx}



%----------------------------------------------------------------------------------------
%	SKILLS
%----------------------------------------------------------------------------------------
\section{Навыки}
\begin{tabularx}{\linewidth}{@{}l X@{}}
ML/DL & \normalsize{Python, PyTorch, Transformers, HuggingFace, scikit-learn, pandas, OpenCV}\\
Инфраструктура & \normalsize{FastAPI, Gradio, Docker, Git}\\
Прочее & \normalsize{C++ (Qt), OpenCL}\\
Личные качества & \normalsize{Обучаемость, умение работать в команде, аналитический склад ума.}\\  
\end{tabularx}

\end{document}
